
\[
\begin{tikzcd}[column sep=large,row sep=large]
\mathord{\wedge} \arrow[r,"\widetilde c"] \arrow[d,"\widetilde c_0"'] &
\Ch(\Gr\mathord{\wedge})\\
\Ch(\Omega) \arrow[ur] &
\end{tikzcd}
\]
D'après la
proposition 6.11, l'homomorphisme : \(\Omega\to\Gr\mathord{\wedge}\) est un isomorphisme si et seulement si l'homomorphisme : \(\Ch(\Omega)\to\Ch(\Gr\mathord{\wedge})\) en est un. Comme il est visiblement surjectif, ce dernier l'est aussi, et il suffit donc de montrer que \(\Ch(\Omega)\to\Ch(\Gr\mathord{\wedge})\) est injectif. Par fonctorialité de \(\Ch\), on obtient le diagramme :
\[
\begin{tikzcd}[column sep=large,row sep=large]
\mathord{\wedge} \arrow[r,"\widetilde c_0"] \arrow[d,"\widetilde c"'] &
\Ch(\Omega) \arrow[d,"\widetilde c"]\\
\Ch(\Gr\mathord{\wedge})) \arrow[r,"\Ch(\widetilde c_0)"] &
\Ch\bigl(\Gr_\lambda(\Ch(\Omega))\bigr);
\end{tikzcd}
\]
comme \(\Omega\) est à degrés bornés, la flèche verticale de droite est un isomorphisme (6.12), et en utilisant 6.11, on obtient donc un morphisme : \(\Gr\mathord{\wedge}\longrightarrow\Omega\), le composé des deux étant l'identité d'après la propriété universelle de \(\Omega\).

\section*{APPENDICE : Les opérations \(\psi^k\) d'Adams}

Nous suivons ici une présentation due à J.P. Serre.

7.1. Soit \(K\) un pré-\(\lambda\)-anneau. Pour tout \(x\in K\), on désigne par \(\frac{d}{dt}(\lambda_t(x))\) la dérivée par rapport à \(t\) de \(\lambda_t(x)\). On définit alors des éléments de \(K\), notés \(\psi^k(x)\), en posant :
\begin{equation}\tag{7.1.1}
  \frac{d}{dt}(\lambda_t(x))/\lambda_t(x)
  =
  \sum_{k=1}^{\infty}(-1)^{k-1}\psi^k(x)t^{k-1}.
\end{equation}

Les applications \(\psi^k:K\to K\) sont additives~; en effet, on a :
\[
\begin{aligned}
 \frac{d}{dt}(\lambda_t(x+y))/\lambda_t(x+y)
 &=
 \frac{d}{dt}(\lambda_t(x)\lambda_t(y))/\lambda_t(x)\lambda_t(y)\\
 &=
 \frac{d}{dt}(\lambda_t(x))/\lambda_t(x)
 +\frac{d}{dt}(\lambda_t(y))/\lambda_t(y)\ ;
\end{aligned}
\]
